% ! TEX root = ../m1-19-20.tex
\chapter{Derivate parțiale}
\index{derivate parțiale}

În general, putem lucra cu funcții de mai multe variabile, sub forma
$ f : \RR^n \to \RR $, funcții care acceptă $ n $ variabile ca \qq{date de intrare}
și rezultatul este o variabilă reală, deci $ f(x_1, \dots, x_n) = y \in \RR $,
pentru orice $ n $-tuplu $ (x_1, \dots, x_n) \in \RR^n $.

Cazurile pe care le vom întîlni cel mai des la acest seminar vor fi, însă,
$ n = 2 $ și $ n = 3 $, cazuri în care variabilele de intrare se vor nota,
respectiv $ (x, y) $ sau $ (x, y, z) $.

Fie, deci, $ f : \RR^3 \to \RR $ o funcție de 3 variabile reale, cu valori
reale, $ f = f(x, y, z) $.

Putem defini \emph{derivata parțială} a funcției $ f $ după variabila $ x $,
de exemplu, ca fiind derivata obținută prin tratarea lui $ y $ și $ z $ ca
parametri (\qq{constante}). Notația este $ \dfrac{\partial f}{\partial x} $
sau $ f_x $ sau $ \partial_x f $.

Similar se pot defini și celelalte derivate parțiale. Un exemplu simplu:
\[
  f : \RR^3 \to \RR, \quad f(x, y, z) = 3x^2 y + 5xe^y + \sin(yz).
\]

Avem:
\begin{align*}
  f_x &= 6xy + 5e^y \\
  f_y &= 3x^2 + 5xe^y + z \cos(yz) \\
  f_z &= y\cos(yz).
\end{align*}

Calculele pot continua și putem defini \emph{derivatele parțiale de ordin %
  superior}:
\[
  f_{xx} = \dfrac{\partial^2 f}{\partial x^2} = %
  \dfrac{\partial}{\partial x} \left( \dfrac{\partial f}{\partial x} \right)
\]
și similar pentru derivate mixte de forma $ f_{xy}, f_{yx} $ sau derivate de ordin
și mai mare.

În cazul funcțiilor pe care le vom folosi, are loc:
\begin{theorem}[Teorema de simetrie a lui Schwarz]\label{thm:schwarz}\index{teorema!Schwarz}
  În anumite ipoteze\footnotemark, derivatele mixte au proprietatea de simetrie,
  adică $ f_{x_i x_j} = f_{x_j x_i} $ sau, scris pe larg:
  \[
    \dfrac{\partial^2 f}{\partial x_i \partial x_j} = \dfrac{\partial^2 f}{\partial x_j \partial x_i},%
    \forall x_i, x_j.
  \]
\end{theorem}
\footnotetext{O prezentare a teoremei se poate găsi \href{https://en.wikipedia.org/wiki/Symmetry\_of\_second\_derivatives\#Schwarz.27s\_theorem}{aici},
  iar în exercițiile de la seminar, ipotezele vor fi verificate automat.}

În exercițiile pe care le vom aborda, vor fi de folos derivatele de ordinul întîi și
cele de ordinul al doilea. În plus, pentru o funcție de 3 variabile reale,
$ f : \RR^3 \to \RR, f = f(x, y, z) $, se definește \emph{laplacianul} (sau
\emph{operatorul Laplace}) prin:
\index{derivate parțiale!laplacian}
\[
  \Delta f = f_{xx} + f_{yy} + f_{zz}.
\]

Similar, desigur, putem defini laplacianul și pentru funcții de 2 variabile reale prin:
\[
  \Delta f = f_{xx} + f_{yy}.
\]

O funcție se numește \emph{armonică} dacă $ \Delta f = 0 $.
\index{derivate parțiale!funcție armonică}

\vspace{1cm}

Dată o funcție de 3 variabile, se poate defini \emph{diferențiala totală},
care conține laolaltă informațiile furnizate de derivatele parțiale. Pentru o
funcție $ f : \RR^3 \to \RR $, diferențiala totală se notează $ df $ și se calculează
cu formula:
\[
  df = f_x dx + f_y dy + f_z dz,
\]
unde $ dx, dy, dz $ sînt diferențialele elementare (nu se calculează! ele
constituie vectori-bază într-un anumit spațiu vectorial).

O expresie ca mai sus, de forma $ df $, se numește \emph{1-formă diferențială},
pe care o vom reîntîlni în studiul integralelor curbilinii.

Avem mai departe și diferențiala totală de ordinul al doilea, definită
pentru aceeași funcție de mai sus prin:
\[
  d^2 f = f_{xx} dx^2 + 2f_{xy} dxdy + f_{yy} dy^2,
\]
expresie care se mai numește \emph{2-formă diferențială} și pe care o
vom reîntîlni în studiul integralelor de suprafață.

\vspace{1cm}

De asemenea, o altă noțiune utilă este aceea a \emph{polinoamelor Taylor}
pentru funcții de mai multe variabile. Fie $ f = f(x, y) $ o funcție
de 2 variabile și fie $ A(x_A, y_A) \in \RR^2 $ un punct arbitrar. În exerciții,
vom întîlni doar polinoamele Taylor de grad 1 și 2, care se definesc astfel,
în jurul punctului $ A $:
\begin{align*}
  T_1(X, Y) &= f(A) + \dfrac{1}{1!}\left[f_x(A)(X - x_A) + f_y(A)(Y - y_A)\right] \\
  T_2(X, Y) &= T_1(X, Y) + \dfrac{1}{2!} \left[ f_{xx}(A)(X - x_A)^2 + f_{yy}(A)(Y - y_A)^2 +%
              2f_{xy}(A)(X - x_A)(Y - y_A) \right].
\end{align*}

Ca în cazul seriilor Taylor pentru funcțiile de o singură variabilă,
eroarea aproximației folosind aceste polinoame este comparabilă cu primul termen neglijat.


%%%%%%%%%%%%%%%%%%%%%%%%%%%%%%%%%%%%%%%%%%%%%%%%%%%%%%%%%%%%%%%%%%%%%%
\section{Extreme libere}
\index{extreme!libere}

Fie o funcție de 2 variabile $ f : \RR^2 \to \RR, f = f(x, y) $. Ne propunem
să studiem valorile sale extreme, adică să îi găsim minimul și maximul,
considerînd că funcția este definită, derivabilă și continuă pe tot domeniul
de definiție.

Pașii pe care îi urmăm sînt:
\begin{enumerate}[(1)]
\item Rezolvăm sistemul de ecuații dat de anularea derivatelor parțiale
  de ordinul întîi, din care aflăm \emph{punctele critice}, care \emph{este posibil}
  să fie de extrem. Așadar, rezolvăm sistemul:
  \[
    \begin{cases}
      f_x &= 0 \\
      f_y &= 0
    \end{cases} \Longrightarrow A_1(x_1, y_1), A_2(x_2, y_2)\dots
  \]
\item Pentru fiecare dintre punctele critice $ A_i $ se alcătuiește
  \emph{matricea hessiană} a funcției, alcătuită din derivatele de ordinul al
  doilea și se evaluează matricea în punctele critice:
  \[
    H_f(A_i) = \begin{pmatrix}
      f_{xx}(A_i) & f_{xy}(A_i) \\
      f_{yx}(A_i) & f_{yy}(A_i)
    \end{pmatrix}
  \]
  \index{derivate parțiale!matricea hessiană}
\item Fie $ A_i $ un punct critic fixat. Calculăm valorile proprii $ \lambda_{1,2} $ ale
  matricei hessiene $ H_f(A_i) $ și decidem astfel:
  \begin{itemize}
  \item Dacă $ \lambda_1, \lambda_2 > 0 $, atunci punctul $ A_i $ este de \emph{minim local};
  \item Dacă $ \lambda_1, \lambda_2 < 0 $, punctul $ A_i $ este \emph{maxim local};
  \item Dacă $ \lambda_1 \cdot \lambda_2 < 0 $, punctul $ A_i $ \emph{nu este de extrem};
  \item Dacă $ \lambda_1 \cdot \lambda_2 = 0 $, nu putem decide.
  \end{itemize}
\item Se repetă procedura pentru fiecare dintre punctele critice.
\end{enumerate}

%%%%%%%%%%%%%%%%%%%%%%%%%%%%%%%%%%%%%%%%%%%%%%%%%%%%%%%%%%%%%%%%%%%%%%
\section{Exerciții}

1. Verificați dacă următoarele funcții de 2 variabile sînt armonice (presupunem
că domeniile de definiție au fost date corect):
\begin{enumerate}[(a)]
\item $ f(x, y) = \ln(x^2 + y^2) $;
\item $ f(x, y) = \sqrt{x^2 + y^2} $;
\item $ f(x, y) = \arctan \dfrac{x}{y} + \arctan\dfrac{y}{x} $;
\item $ f(x, y) = \dfrac{x + y}{x - y} $.
\end{enumerate}

\vspace{1cm}

2. Găsiți punctele de extrem pentru funcțiile de 2 sau 3 variabile, definite
corespunzător pe $ \RR^2 $ sau $ \RR^3 $:
\begin{enumerate}[(a)]
\item $ f(x, y) = 3x^2y - 2x + 5y - xy^2 $;
\item $ f(x, y) = x^3 + y^3 - 6xy $;
\item $ f(x, y) = 3xy^2 - x^3 - 15x - 36y + 9 $;
\item $ f(x, y) = 4xy - x^4 - y^4 $;
\end{enumerate}

%%% Local Variables:
%%% mode: latex
%%% TeX-master: "../m1-19-20"
%%% End:
